\scanpage{6}
\mspage{\handwritten{(29}}
\begingroup
\makeatletter\tagsleft@true\makeatother
\[
\underline{\mathrm{Pt}}(X)^{\unsure{\circ|o|0}}\dashrightarrow
\underline{\mathrm{Pt}}(\handwritten{\widetilde}{X_{\text{ét}}})\ ,
\tag*{$(*)$}
\]
where the left-hand side is defined in 7.6.. This being so, we have the

\noindent\uline{Theorem} 7.9.\quad Let $X$ be a prescheme. The preceding functor $(*)$ is an \uline{an\deleted{\unsure{it|ti}}\handwritten{ti}-equivalence} of the category of geometric points over $X$ (with the specialization arrows as morphisms) with the category of fiber functors on the étale topos $\handwritten{\widetilde}{X_{\text{ét}}}$.

As this theorem will not be used again in the rest of the seminar we confine ourselves here to a sketch of proof, where we shall allow ourselves certain liberties with questions of Universes.

a) \uline{The functor under consideration is fully faithful}. For every $X\handwritten{\in}\operatorname{Ob} X_{\text{ét}}$, let $X^{\handwritten{\unsure{\vee|v}}}: \handwritten{\widetilde}{X_{\text{ét}}} \dashrightarrow (\mathrm{Ens})$ be defined by
\[
\handwritten{\check}{X}(F) = \operatorname{Hom}(X,F) = F(X).
\]

\noindent Note that the fiber functor
\[
\handwritten{\varepsilon_{\xi}} : F \handwritten{\rightsquigarrow} F_{\handwritten{\xi}}
\]
can be written
\[
\handwritten{\varepsilon_{\xi}} \handwritten{\simeq} \mathop{\underset{\handwritten{\longrightarrow}}{\mathrm{lim}}}\limits_i X_i^{\handwritten{\vee}}\ ,
\]
where the $X_i$ are affine preschemes étale over $X$ indexed by a certain pseudo-filtering category (cf 4.3. and 4.5.). We therefore have for every functor $\handwritten{\varphi}: \handwritten{\widetilde}{X_{\text{ét}}} \dashrightarrow (\mathrm{Ens})$:\margin{\tikz{\fill (0,0) circle (.018);\fill (.08,.12) circle (.018);}}
\[
\operatorname{Hom}(\handwritten{\varepsilon_{\xi}},\handwritten{\varphi}) \mathrel{\mathop{\text{-\kern.04em-}}\limits^{\handwritten{\sim}}} \mathop{\underset{\handwritten{\longleftarrow}}{\mathrm{lim}}}\limits_{\handwritten{i}} \operatorname{Hom}(X_i^{\handwritten{\vee}},\handwritten{\varphi})\ ,
\]
on the other hand it is well known that we have a bifunctorial isomorphism
\[
\operatorname{Hom}(X_i^{\handwritten{\vee}},\handwritten{\varphi}) \mathrel{\mathop{\text{-\kern.04em-}}\limits^{\handwritten{\sim}}} \handwritten{\varphi}(X_i)\ .
\]

\figure{\begin{tikzpicture}[x=1cm,y=1cm]
\draw (0,0) .. controls (.45,-.76) and (1.34,-1.79) .. (1.83,-2.47);
\end{tikzpicture}}
\endgroup

% archive sheet number: 5
